\section{Data Dictionary}\label{sec:dict}

The tables below reproduce the \emph{ACTUS Data Dictionary} (ACTUS Financial Research Foundation, v1.4, CC-BY-SA-4.0), the normative glossary of the standard.  They explain every contract-term, state-variable and contract-type acronym used throughout this document; the lifecycle event types (\texttt{IED}, \texttt{IP}, \texttt{PR}, \dots) are not part of the dictionary and so are not listed here.  Acronyms in \texttt{typewriter} font elsewhere in the blueprint link to their entry below.

\subsection{Contract Terms}

\begin{dicttable}
	\hypertarget{term:AMD}{\texttt{AMD}} & Amortization Date & This Date is used to calculate the annuity amounts for ANN and ANX NGX CT's. Needs only to be set in case where the contract balloon at MD and MD is less than AD. \\ \hline
	\hypertarget{term:ARFIXVAR}{\texttt{ARFIXVAR}} & Array Fixed Variable & For array-type rate reset schedules, this attributes defines the meaning of ARRATE. \\ \hline
	\hypertarget{term:ARINCDEC}{\texttt{ARINCDEC}} & Array Increase Decrease & Indicates whether a certain PRNXT element in ARPRNX increases the principal (NT) or decreases it. Applies only for ANX, NAX, LAX Maturity CTs. For all other Maturity CTs the first principal payment is always in the opposite direction of all other (following) payments. \\ \hline
	\hypertarget{term:ARIPANXi}{\texttt{ARIPANXi}} & Array Cycle Anchor Date Of Interest Payment & Same like IPANX but as array \\ \hline
	\hypertarget{term:ARIPCLi}{\texttt{ARIPCLi}} & Array Cycle Of Interest Payment & Same like IPCL but as array \\ \hline
	\hypertarget{term:ARPRANXj}{\texttt{ARPRANXj}} & Array Cycle Anchor Date Of Principal Redemption & Same like PRANX but as array \\ \hline
	\hypertarget{term:ARPRCLj}{\texttt{ARPRCLj}} & Array Cycle Of Principal Redemption & Same like PRCL but as array \\ \hline
	\hypertarget{term:ARPRNXTj}{\texttt{ARPRNXTj}} & Array Next Principal Redemption Payment & Same like PRNXT but as array \\ \hline
	\hypertarget{term:ARRATE}{\texttt{ARRATE}} & Array Rate & For array-type rate reset schedules, this attribute represents either an interest rate (corresponding to IPNR) or a spread (corresponding to RRSP). Which case applies depends on the attribute ARFIXVAR: if ARFIXVAR=FIX then it represents the new IPNR and if ARFIXVAR=VAR then the applicable RRSP. \\ \hline
	\hypertarget{term:ARRRANX}{\texttt{ARRRANX}} & Array Cycle Anchor Date Of Rate Reset & Same like RRANX but as array \\ \hline
	\hypertarget{term:ARRRCL}{\texttt{ARRRCL}} & Array Cycle Of Rate Reset & Same like RRCL but as array \\ \hline
	\hypertarget{term:BCF}{\texttt{BCF}} & Boundary Crossed Flag & Initializes the value of Boundary Crossed Flag state variable at statusDate \\ \hline
	\hypertarget{term:BDC}{\texttt{BDC}} & Business Day Convention & BDC's are linked to a calendar. Calendars have working and non-working days. A BDC value other than N means that cash flows cannot fall on non-working days, they must be shifted to the next business day (following) or the previous on (preceding). These two simple rules get refined twofold: - Following modified (preceding): Same like following (preceding), however if a cash flow gets shifted into a new month, then it is shifted to preceding (following) business day. - Shift/calculate (SC) and calculate/shift (CS). Accrual, principal, and possibly other calculations are affected by this choice. In the case of SC first the dates are shifted and after the shift cash flows are calculated. In the case of CS it is the other way round. Attention: Does not affect non-cyclical dates such as PRD, MD, TD, IPCED since they can be set to the correct date directly. \\ \hline
	\hypertarget{term:BDR}{\texttt{BDR}} & Boundary Direction & Boundary direction specifies the direction of motion in the underlying asset's price which will be considered a valid crossing of the boundary and trigger the boundary effect changing which, if any, of the boundary legs is active. \\ \hline
	\hypertarget{term:BEF}{\texttt{BEF}} & Boundary Effect & This term specifies which leg - if any- becomes the active subcontract when the underlying asset's price crosses the specified boundary value in the specified direction triggerring a boundary crossing event. \\ \hline
	\hypertarget{term:BLIA}{\texttt{BLIA}} & Boundary Leg Initially Active & Specifies which leg - if any - is the active contract in effect when the boundary controlled switch contract starts. \\ \hline
	\hypertarget{term:BMANX}{\texttt{BMANX}} & Boundary Monitoring Anchor Date & The first Boundary monitoring event occurs on this date \\ \hline
	\hypertarget{term:BMCL}{\texttt{BMCL}} & Boundary Monitoring Cycle & The frequency with which boundary monitoring events occur. It defines how often the system checks to test whether the market value of the underlying asset has crossed the boundary in the specified direction triggerring a boundary crossing event. \\ \hline
	\hypertarget{term:BMED}{\texttt{BMED}} & Boundary Monitoring End Date & Boundary monitoring ends on this date \\ \hline
	\hypertarget{term:BV}{\texttt{BV}} & Boundary Value & Boundary value in a barrier options contract, when reached, triggers the boundary effect specified e.g. Knock-In or Knock-out \\ \hline
	\hypertarget{term:CDD}{\texttt{CDD}} & Contract Deal Date & This date signifies the origination of the contract where an agreement between the customer and the bank has been settled. From this date on, the institution will have a (market) risk position for financial contracts. This is even the case when IED is in future. \\ \hline
	\hypertarget{term:CECV}{\texttt{CECV}} & Coverage Of Credit Enhancement & Defines which percentage of the exposure is covered \\ \hline
	\hypertarget{term:CEGE}{\texttt{CEGE}} & Guaranteed Exposure & Defines which value of the exposure is covered: - NO: Nominal Value - NI: Nominal plus Interest - MV: Market Value \\ \hline
	\hypertarget{term:CETC}{\texttt{CETC}} & Credit Event Type Covered & The type of credit events covered e.g. in credit enhancement or credit default swap contracts. Only the defined credit event types may trigger the protection. \\ \hline
	\hypertarget{term:CID}{\texttt{CID}} & Contract Identifier & Unique identifier of a contract. If the system is used on a single firm level, an internal unique ID can be generated. If used on a national or globally level, a globally unique ID is required. \\ \hline
	\hypertarget{term:CLA}{\texttt{CLA}} & Credit Line Amount & If defined, gives the total amount that can be drawn from a credit line. The remaining amount that can still be drawn is given by CLA-NT. For ANN, NAM, the credit line can only be drawn prior to PRANX-1PRCL. For CRL, the remaining amount that can still be drawn is given by CLA-Sum(NT of attached contracts). \\ \hline
	\hypertarget{term:CLDR}{\texttt{CLDR}} & Calendar & Calendar defines the non-working days which affect the dates of contract events (CDE's) in combination with EOMC and BDC. Custom calendars can be added as additional enum options. \\ \hline
	\hypertarget{term:CNTRL}{\texttt{CNTRL}} & Contract Role & CNTRL defines which position the CRID ( the creator of the contract record ) takes in a contract. For example, whether the contract is an asset or liability, a long or short position for the CRID. Most contracts are simple on or off balance sheet positions which are assets, liabilities. Such contracts can also play a secondary role as a collateral. The attribute is highly significant since it determines the direction of all cash flows. The exact meaning is given with each CT in the ACTUS High Level Specification document. \\ \hline
	\hypertarget{term:CPID}{\texttt{CPID}} & Counterparty Identifier & CPID identifies the counterparty to the CRID in this contract. CPID is ideally the official LEI which can be a firm, a government body, even a single person etc. However, this can also refer to a annonymous group in which case this information is not to be disclosed. CPID may also refer to a group taking a joint risk or more generally, CPID is the main counterparty, against which the contract has been settled. \\ \hline
	\hypertarget{term:CRID}{\texttt{CRID}} & Creator Identifier & This identifies the legal entity creating the contract record. The counterparty of the contract is tracked in CPID. CRID is ideally the official LEI which can be a firm, a government body, even a single person etc. However, this can also refer to a annonymous group in which case this information is not to be disclosed. CRID may also refer to a group taking a joint risk. \\ \hline
	\hypertarget{term:CT}{\texttt{CT}} & Contract Type & The ContractType is the most important information. It defines the cash flow generating pattern of a contract. The ContractType information in combination with a given state of the risk factors will produce a deterministic sequence of cash flows which are the basis of any financial analysis. \\ \hline
	\hypertarget{term:CTS}{\texttt{CTS}} & Contract Structure & A structure identifying individual or sets of underlying contracts. E.g. for FUTUR, this structure identifies the single underlying contract, for SWAPS, the FirstLeg and SecondLeg are identified, or for CEG, CEC the structure identifies Covered and Covering contracts. \\ \hline
	\hypertarget{term:CUR}{\texttt{CUR}} & Currency & The currency of the cash flows. \\ \hline
	\hypertarget{term:CUR2}{\texttt{CUR2}} & Currency 2 & The currency of the cash flows of the second leg (if not defined, main currency applies) \\ \hline
	\hypertarget{term:CURS}{\texttt{CURS}} & Settlement Currency & The currency in which cash flows are settled. This currency can be different from the currency (CUR) in which cash flows or the contract, respectively, is denominated in which case the respective FX-rate applies at settlement time. If no settlement currency is defined the cash flows are settled in the currency in which they are denominated. \\ \hline
	\hypertarget{term:DQP}{\texttt{DQP}} & Delinquency Period & If real payment happens after scheduled payment date plus DLP, then the counterparty is in technical default. This means that the creditor legally has the right to declare default of the debtor. \\ \hline
	\hypertarget{term:DQR}{\texttt{DQR}} & Delinquency Rate & Rate at which Delinquency Payments accrue on NT (in addition to the interest rate) during the DelinquencyPeriod \\ \hline
	\hypertarget{term:DS}{\texttt{DS}} & Delivery Settlement & Indicates whether the contract is settled in cash or physical delivery. In case of physical delivery, the underlying contract and associated (future) cash flows are effectively exchanged. In case of cash settlement, the current market value of the underlying contract determines the cash flow exchanged. \\ \hline
	\hypertarget{term:DVANX}{\texttt{DVANX}} & Cycle Anchor Date Of Dividend & Date from which the dividend payment date schedule is calculated according to the cycle length. The first dividend payment event takes place on this anchor. \\ \hline
	\hypertarget{term:DVCL}{\texttt{DVCL}} & Cycle Of Dividend & Defines in combination with DVANX the payment points of dividends. The dividend payment schedule will start at DVANX and end at MaximumProjectionPeriod (cf. sheet Modeling Parameters). \\ \hline
	\hypertarget{term:DVEX}{\texttt{DVEX}} & Ex Dividend Date & In case contract is traded between DVEX and next DV payment date (i.e. PRD>DVEX \& PRD<next DV payment date), then the old holder of the contract (previous to the trade) receives the next DV payment. In other words, the next DV payment is cancelled for the new (after the trade) holder of the contract. \\ \hline
	\hypertarget{term:DVNP}{\texttt{DVNP}} & Next Dividend Payment Amount & Defines the next dividend payment (amount) whereas the date of dividend payment is defined through the DVANX/DVCL pair. If DVCL is defined, then this amount will be used as dividend payment for each future dividend payment date. \\ \hline
	\hypertarget{term:EOMC}{\texttt{EOMC}} & End Of Month Convention & When computing schedules a special problem arises if an anchor date is at the end of a month and a cycle of monthly or quarterly is applied (yearly in the case of leap years only). How do we have to interpret an anchor date April 30 plus 1M cycles? In case where EOM is selected, it will jump to the 31st of May, then June 30, July 31 and so on. If SM is selected, it will jump to the 30st always with of course an exception in February. This logic applies for all months having 30 or less days and an anchor date at the last day. Month with 31 days will at any rate jump to the last of the month if anchor date is on the last day. \\ \hline
	\hypertarget{term:FEAC}{\texttt{FEAC}} & Fee Accrued & Accrued fees as per SD \\ \hline
	\hypertarget{term:FEANX}{\texttt{FEANX}} & Cycle Anchor Date Of Fee & Date from which the fee payment date schedule is calculated according to the cycle length. The first fee payment event takes place on this anchor. \\ \hline
	\hypertarget{term:FEB}{\texttt{FEB}} & Fee Basis & Basis, on which Fee is calculated. For FEB=’A’, FER is interpreted as an absolute amount to be paid at every FP event and for FEB=’N’, FER represents a rate at which FP amounts accrue on the basis of the contract’s NT. \\ \hline
	\hypertarget{term:FECL}{\texttt{FECL}} & Cycle Of Fee & Defines in combination with FEANX the payment points of fees \\ \hline
	\hypertarget{term:FER}{\texttt{FER}} & Fee Rate & Rate of Fee which is a percentage of the underlying or FER is an absolute amount. For all contracts where FEB does not apply (cf. business rules), FER is interpreted as an absolute amount. \\ \hline
	\hypertarget{term:GRP}{\texttt{GRP}} & Grace Period & If real payment happens after scheduled payment date plus GRP, then the payment is in delay. \\ \hline
	\hypertarget{term:IED}{\texttt{IED}} & Initial Exchange Date & Date of the initial cash flow for Maturity and Non-Maturity CT's. It also coincides with the beginning of interest accrual calculation. \\ \hline
	\hypertarget{term:IPAC}{\texttt{IPAC}} & Accrued Interest & Accrued interest as per SD. In case of NULL, this value will be recalculated using IPANX, IPCL and IPNR information. Can be used to represent irregular next IP payments. \\ \hline
	\hypertarget{term:IPANX}{\texttt{IPANX}} & Cycle Anchor Date Of Interest Payment & Date from which the interest payment date schedule is calculated according to the cycle length. The first interest payment event takes place on this anchor. \\ \hline
	\hypertarget{term:IPCB}{\texttt{IPCB}} & Interest Calculation Base & This is important for amortizing instruments. The basis of interest calculation is normally the notional outstanding amount as per SD. This is considered the fair basis and in many countries the only legal basis. If NULL or NTSD is selected, this is the case. Alternative bases (normally in order to favor the lending institution) are found. In the extreme case the original balance (PCDD=NT+PDCDD) never gets adjusted. In this case PCDD must be chosen. An intermediate case exist wherre balances do get adjusted, however with lags. In this case NTL mut be selected and anchor dates and cycles must be set. \\ \hline
	\hypertarget{term:IPCBA}{\texttt{IPCBA}} & Interest Calculation Base Amount & This is the amount used for the calculation of interest. Calculation base per SD. \\ \hline
	\hypertarget{term:IPCBANX}{\texttt{IPCBANX}} & Cycle Anchor Date Of Interest Calculation Base & Date from which the interest calculation base date schedule is calculated according to the cycle length. The first interest calculation base event takes place on this anchor. \\ \hline
	\hypertarget{term:IPCBCL}{\texttt{IPCBCL}} & Cycle Of Interest Calculation Base & Concerning the format see PRCL. Defines the subsequent adjustment points to NT of the interest payment calculation base. \\ \hline
	\hypertarget{term:IPCED}{\texttt{IPCED}} & Capitalization End Date & If IPCED is set, then interest is not paid or received but added to the balance (NT) until IPCED. If IPCED does not coincide with an IP cycle, one additional interest payment gets calculated at IPCED and capitalized. Thereafter normal interest payments occur. \\ \hline
	\hypertarget{term:IPCL}{\texttt{IPCL}} & Cycle Of Interest Payment & Cycle according to which the interest payment date schedule will be calculated. In case IPCL is not set, then there will only be an interest payment event at MD (and possibly at IPANX if set). The interval will be adjusted yet by EOMC and BDC. \\ \hline
	\hypertarget{term:IPDC}{\texttt{IPDC}} & Day Count Convention & Method defining how days are counted between two dates. This finally defines the year fraction in accrual calculations. \\ \hline
	\hypertarget{term:IPNR}{\texttt{IPNR}} & Nominal Interest Rate & The nominal interest rate which will be used to calculate accruals and the next interest payment at the next IP date. NT multiplied with IPNR is the base for the interest payment calculation. The relevant time period is a function of IPDC. If the contract is variable (RRANX set) this field is periodically updated per SD. In the case of plan vanilla interest rate swaps (IRSPV) this defines the rate of fixed leg. \\ \hline
	\hypertarget{term:IPNR2}{\texttt{IPNR2}} & Nominal Interest Rate 2 & The nominal interest rate which will be used to calculate accruals and the next interest payment at the next IP date on the second leg (the one not mentioned in CNTRL) of a plain vanilla swap. The relevant time period is a function of IPDC. It is periodically updated per SD. \\ \hline
	\hypertarget{term:IPPNT}{\texttt{IPPNT}} & Cycle Point Of Interest Payment & Usually, interest is paid at the end of each IPCL which corresponds to a IPPNT value of E which is also the default. If interest payment occurs at the beginning of the cycle, the value is B. \\ \hline
	\hypertarget{term:MD}{\texttt{MD}} & Maturity Date & Marks the contractual end of the lifecycle of a CT. Generally, date of the last cash flows. This includes normally a principal and an interest payment. Some Maturity CTs as perpetuals (PBN) do not have such a date. For variable amortizing contracts of the ANN CT, this date might be less than the scheduled end of the contract (which is deduced from the periodic payment amount PRNXT). In this case it balloons. \\ \hline
	\hypertarget{term:MOC}{\texttt{MOC}} & Market Object Code & Is pointing to the market value at SD (MarketObject). Unique codes for market objects must be used. \\ \hline
	\hypertarget{term:MPFD}{\texttt{MPFD}} & Maximum Penalty Free Disbursement & Defines the notional amount which can be withdrawn before XDN without penalty \\ \hline
	\hypertarget{term:MRANX}{\texttt{MRANX}} & Cycle Anchor Date Of Margining & Date from which the margin call date schedule is calculated according to the cycle length. The first margin call event takes place on this anchor. \\ \hline
	\hypertarget{term:MRCL}{\texttt{MRCL}} & Cycle Of Margining & Defines together with MRANX the points where margins can be called. \\ \hline
	\hypertarget{term:MRCLH}{\texttt{MRCLH}} & Clearing House & Indicates wheter CRID takes a clearing house function or not. In other word, whether CRID receive margins (MRIM, MRVM). \\ \hline
	\hypertarget{term:MRIM}{\texttt{MRIM}} & Initial Margin & Margin to cover losses which may be incurred as a result of market fluctuations. Upon contract closing or maturity, the MRIM is reimbursed. \\ \hline
	\hypertarget{term:MRMML}{\texttt{MRMML}} & Maintenance Margin Lower Bound & Defines the lower bound of the Maintenance Margin. If MRVM falls below MRMML, then capital must be added to reach the original MRIM. \\ \hline
	\hypertarget{term:MRMMU}{\texttt{MRMMU}} & Maintenance Margin Upper Bound & Defines the upper bound of the Maintenance Margin. If MRVM falls above MRMMU, then capital is refunded to reach the original MRIM. \\ \hline
	\hypertarget{term:MRVM}{\texttt{MRVM}} & Variation Margin & MRVM reflects the accrued but not yet paid margin as per SD. Open traded positions are revalued by the exchange at the end of every trading day using mark-to-market valuation. Often clearing members do not credit or debit their clients daily with MRVM, but rather use a Maintenance Margin. If the balance falls outside MRMML (and MRMMU), then capital must be added (is refunded) to reach the original margin amount MRIM. We can also say that MVO+MRVM is equal to the reference value as per last margin update. \\ \hline
	\hypertarget{term:MVO}{\texttt{MVO}} & Market Value Observed & Value as observed in the market at SD per unit. Incase of fixed income instruments it is a fraction. \\ \hline
	\hypertarget{term:NPD}{\texttt{NPD}} & Non Performing Date & The date of the (uncovered) payment event responsible for the current value of the Contract Performance attribute. \\ \hline
	\hypertarget{term:NT}{\texttt{NT}} & Notional Principal & Current nominal value of the contract. For debt instrument this is the current remaining notional outstanding. NT is generally the basis on which interest payments are calculated. If IPCBS is set, IPCBS may introduce a different basis for interest payment calculation. \\ \hline
	\hypertarget{term:NT2}{\texttt{NT2}} & Notional Principal 2 & Notional amount of the second currency to be exchanged in an FXOUT CT. \\ \hline
	\hypertarget{term:OPANX}{\texttt{OPANX}} & Cycle Anchor Date Of Optionality & Used for Basic Maturities (such as PAM, RGM, ANN, NGM and their Step-up versions) and American and Bermudan style options. - Basic Maturities: Within the group of these Maturities, it indicates the possibility of prepayments. Prepayment features are controlled by Behavior. - American and Bermudan style Options: Begin of exercise period. \\ \hline
	\hypertarget{term:OPCL}{\texttt{OPCL}} & Cycle Of Optionality & Cycle according to which the option exercise date schedule will be calculated. OPCL can be NULL for American Options or Prepayment Optionality in which case the optionality period starts at OPANX and ends at OPXED (for american options) or MD (in case of prepayment optionality). The interval will be adjusted yet by EOMC and BDC. \\ \hline
	\hypertarget{term:OPS1}{\texttt{OPS1}} & Option Strike 1 & Strike price of the option. Whether it is a call/put is determined by the attribute OPTP, i.e a call or a put (or a combination of call/put). This attribute is used for price related options such as options on bonds, stocks or FX. Interest rate related options (caps/floos) are handled within th RatReset group. \\ \hline
	\hypertarget{term:OPTP}{\texttt{OPTP}} & Option Type & Defines whether the option is a call or put or a combination of it. This field has to be seen in combination with CNTRL where it is defined whether CRID is the buyer or the seller. \\ \hline
	\hypertarget{term:OPXED}{\texttt{OPXED}} & Option Exercise End Date & Final exercise date for American and Bermudan options, expiry date for European options. \\ \hline
	\hypertarget{term:OPXT}{\texttt{OPXT}} & Option Exercise Type & Defines whether the option is European (exercised at a specific date), American (exercised during a span of time) or Bermudan (exercised at certain points during a span of time). \\ \hline
	\hypertarget{term:PDIED}{\texttt{PDIED}} & Premium Discount At IED & Total original premium or discount that has been set at CDD and will be added to the (notional) cash flow at IED (cash flow at IED = NT+PDIED, w.r.t. an RPA CT). Negative value for discount and positive for premium. Note, similar to interest the PDIED portion is part of P\&L. \\ \hline
	\hypertarget{term:PFUT}{\texttt{PFUT}} & Futures Price & The price the counterparties agreed upon at which the underlying contract (of a FUTUR) is exchanged/settled at STD. Quoting is different for different types of underlyings: Fixed Income = in percentage, all others in nominal terms. \\ \hline
	\hypertarget{term:PPEF}{\texttt{PPEF}} & Prepayment Effect & This attribute defines whether or not the right of prepayment exists and if yes, how prepayment affects the remaining principal redemption schedule of the contract \\ \hline
	\hypertarget{term:PPP}{\texttt{PPP}} & Prepayment Period & If real payment happens before scheduled payment date minus PPP, then it is considered a prepayment. Effect of prepayments are further described in PPEF and related fields. \\ \hline
	\hypertarget{term:PPRD}{\texttt{PPRD}} & Price At Purchase Date & Purchase price exchanged at PRD. PPRD represents a clean price (includes premium/discount but not IPAC). \\ \hline
	\hypertarget{term:PRANX}{\texttt{PRANX}} & Cycle Anchor Date Of Principal Redemption & Date from which the principal payment date schedule is calculated according to the cycle length. The first principal payment event takes place on this anchor. \\ \hline
	\hypertarget{term:PRCL}{\texttt{PRCL}} & Cycle Of Principal Redemption & Cycle according to which the interest payment date schedule will be calculated. In case PRCL is not set, then there will only be one principal payment event at MD (and possibly at PRANX if set). The interval will be adjusted yet by EOMC and BDC. \\ \hline
	\hypertarget{term:PRD}{\texttt{PRD}} & Purchase Date & If a contract is bought after initiation (for example a bond on the secondary market) this date has to be set. It refers to the date at which the payment (of PPRD) and transfer of the security happens. In other words, PRD - if set - takes the role otherwise IED has from a cash flow perspective. Note, CPID of the CT is not the counterparty of the transaction! \\ \hline
	\hypertarget{term:PRF}{\texttt{PRF}} & Contract Performance & Indicates the current contract performance status. Different states of the contract range from performing to default. \\ \hline
	\hypertarget{term:PRNXT}{\texttt{PRNXT}} & Next Principal Redemption Payment & Amount of principal that will be paid during the redemption cycle at the next payment date. For amortizing contracts like ANN, NAM, ANX, and NAX this is the total periodic payment amount (sum of interest and principal). \\ \hline
	\hypertarget{term:PTD}{\texttt{PTD}} & Price At Termination Date & Sellingprice exchanged at PTD PTDrepresents a clean price (includes premium/discount but not IPAC \\ \hline
	\hypertarget{term:PYRT}{\texttt{PYRT}} & Penalty Rate & Either the rate or the absolute amount of the prepayment. \\ \hline
	\hypertarget{term:PYTP}{\texttt{PYTP}} & Penalty Type & Defines whether prepayment is linked to a penalty and of which kind. \\ \hline
	\hypertarget{term:QT}{\texttt{QT}} & Quantity & This attribute relates either to physical contracts (COM) or underlyings of traded contracts. In case of physical contracts it holds the number of underlying units of the specific good (e.g. number of barrels of oil). In case of well defined traded contracts it holds the number of defined underlying instruments. Example: QT of STK CTs underlying a FUTUR indicates the number of those specific STK CTs which underlie the FUTUR. \\ \hline
	\hypertarget{term:RRANX}{\texttt{RRANX}} & Cycle Anchor Date Of Rate Reset & Date from which the rate reset date schedule is calculated according to the cycle length. The first rate reset event takes place on this anchor. \\ \hline
	\hypertarget{term:RRCL}{\texttt{RRCL}} & Cycle Of Rate Reset & Cycle according to which the rate reset date schedule will be calculated. In case RRCL is not set, then there will only be one rate reset event at RRANX given RRANX if set. The interval will be adjusted yet by EOMC and BDC. \\ \hline
	\hypertarget{term:RRFIX}{\texttt{RRFIX}} & Fixing Period & Interest rate resets (adjustments) are usually fixed one or two days (usually Business Days) before the new rate applies (defined by the rate reset schedule). This field holds the period between fixing and application of a rate. \\ \hline
	\hypertarget{term:RRLC}{\texttt{RRLC}} & Life Cap & For variable rate basic CTs this represents a cap on the interest rate that applies during the entire lifetime of the contract. For CAPFL CTs this represents the cap strike rate. \\ \hline
	\hypertarget{term:RRLF}{\texttt{RRLF}} & Life Floor & For variable rate basic CTs this represents a floor on the interest rate that applies during the entire lifetime of the contract. For CAPFL CTs this represents the floor strike rate. \\ \hline
	\hypertarget{term:RRMLT}{\texttt{RRMLT}} & Rate Multiplier & Interest rate multiplier. A typical rate resetting rule is LIBOR plus x basis point where x represents the interest rate spread. However, in some cases like reverse or super floater contracts an additional rate multiplier applies. In this case, the new rate is determined as: IPNR after rate reset = Rate selected from the market object * RRMLT + RRSP. \\ \hline
	\hypertarget{term:RRMO}{\texttt{RRMO}} & Market Object Code Of Rate Reset & Is pointing to the interest rate driver (MarketObject) used for rate reset uniquely. Unique codes for market objects must be used. \\ \hline
	\hypertarget{term:RRNXT}{\texttt{RRNXT}} & Next Reset Rate & Holds the new rate that has been fixed already (cf. attribute FixingDays) but not applied. This new rate will be applied at the next rate reset event (after SD and according to the rate reset schedule). Attention, RRNXT must be set to NULL after it is applied! \\ \hline
	\hypertarget{term:RRPC}{\texttt{RRPC}} & Period Cap & For variable rate basic CTs this represents the maximum positive rate change per rate reset cycle. \\ \hline
	\hypertarget{term:RRPF}{\texttt{RRPF}} & Period Floor & For variable rate basic CTs this represents the maximum negative rate change per rate reset cycle. \\ \hline
	\hypertarget{term:RRPNT}{\texttt{RRPNT}} & Cycle Point Of Rate Reset & Normally rates get reset at the beginning of any resetting cycles. There are contracts where the rate is not set at the beginning but at the end of the cycle and then applied to the previous cycle (post-fixing); in other words the rate applies before it is fixed. Hence, the new rate is not known during the entire cycle where it applies. Therefore, the rate will be applied backwards at the end of the cycle. This happens through a correction of interest accrued. \\ \hline
	\hypertarget{term:RRSP}{\texttt{RRSP}} & Rate Spread & Interest rate spread. A typical rate resetting rule is LIBOR plus x basis point where x represents the interest rate spread. The following equation can be taken if RRMLT is not set: IPNR after rate reset = Rate selected from the market object + RRSP. \\ \hline
	\hypertarget{term:SCANX}{\texttt{SCANX}} & Cycle Anchor Date Of Scaling Index & Date from which the scaling date schedule is calculated according to the cycle length. The first scaling event takes place on this anchor. \\ \hline
	\hypertarget{term:SCCDD}{\texttt{SCCDD}} & Scaling Index At Contract Deal Date & The value of the Scaling Index as per Contract Deal Date. \\ \hline
	\hypertarget{term:SCCL}{\texttt{SCCL}} & Cycle Of Scaling Index & Cycle according to which the scaling date schedule will be calculated. In case SCCL is not set, then there will only be one scaling event at SCANX given SCANX is set. The interval will be adjusted yet by EOMC and BDC. \\ \hline
	\hypertarget{term:SCEF}{\texttt{SCEF}} & Scaling Effect & Indicates which payments are scaled. I = Interest payments, N = Nominal payments and M = Maximum deferred interest amount. They can be scaled in any combination. \\ \hline
	\hypertarget{term:SCIP}{\texttt{SCIP}} & Interest Scaling Multiplier & The multiplier being applied to interest cash flows \\ \hline
	\hypertarget{term:SCMO}{\texttt{SCMO}} & Market Object Code Of Scaling Index &  \\ \hline
	\hypertarget{term:SCNT}{\texttt{SCNT}} & Notional Scaling Multiplier & The multiplier being applied to principal cash flows \\ \hline
	\hypertarget{term:SD}{\texttt{SD}} & Status Date & SD holds the date per which all attributes of the record were updated. This is especially important for the highly dynamic attributes like Accruals, Notional, interest rates in variable instruments etc. \\ \hline
	\hypertarget{term:SEN}{\texttt{SEN}} & Seniority & Refers to the order of repayment in the event of a sale or default of the issuer. \\ \hline
	\hypertarget{term:STP}{\texttt{STP}} & Settlement Period & Defines the period from fixing of a contingent event/obligation (Exercise Date) to settlement of the obligation. \\ \hline
	\hypertarget{term:TD}{\texttt{TD}} & Termination Date & If a contract is sold before MD (for example a bond on the secondary market) this date has to be set. It refers to the date at which the payment (of PTD) and transfer of the security happens. In other words, TD - if set - takes the role otherwise MD has from a cash flow perspective. Note, CPID of the CT is not the counterparty of the transaction! \\ \hline
	\hypertarget{term:UT}{\texttt{UT}} & Unit & The physical unit of the contract. Example: Barrels for an Oil COM CT. \\ \hline
	\hypertarget{term:XA}{\texttt{XA}} & Exercise Amount & The amount fixed at Exercise Date for a contingent event/obligation such as a forward condition, optionality etc. The Exercise Amount is fixed at Exercise Date but not settled yet. \\ \hline
	\hypertarget{term:XD}{\texttt{XD}} & Exercise Date & Date of exercising a contingent event/obligation such as a forward condition, optionality etc. The Exercise date marks the observed timestamp of fixing the contingent event and respective payment obligation not necessarily the timestamp of settling the obligation. \\ \hline
	\hypertarget{term:XDN}{\texttt{XDN}} & X Day Notice & Used as rolling attribute with the callable CT's UMP and CLM uniquely. CLM's and UMP's will not be settled (MD not set) until the client uses his option to call the contract X\_Day\_Notice after Current Date. As long as MD or TD is not set, the client postpones his right to call to the future. The cycle is normally defined in number of business days. \\ \hline
	\hypertarget{term:sss}{\texttt{sss}} & Option Strike 2 & Put price in case of call/put. \\ \hline
\end{dicttable}

\subsection{State Variables}

\begin{dicttable}
	\hypertarget{state:Bcf}{\texttt{Bcf}} & Boundary Crossed Flag & TRUE value indicates that a boundary crossing event has already occurred; boundary monitoring is no longer in effect. \\ \hline
	\hypertarget{state:Blaf1}{\texttt{Blaf1}} & Boundary Leg1 Active Flag & TRUE value indicates that boundaryLeg1 is currently active and generating cashflow events. \\ \hline
	\hypertarget{state:Blaf2}{\texttt{Blaf2}} & Boundary Leg2 Active Flag & TRUE value indicates that boundaryLeg2 is currently active and generating cashflow events. \\ \hline
	\hypertarget{state:Feac}{\texttt{Feac}} & Fee Accrued & The current value of accrued fees \\ \hline
	\hypertarget{state:Icba}{\texttt{Icba}} & Interest Calculation Base Amount & The basis at which interest is being accrued. Potentially different from NVL. \\ \hline
	\hypertarget{state:Ipac}{\texttt{Ipac}} & Accrued Interest & The current value of accrued interest \\ \hline
	\hypertarget{state:Ipac2}{\texttt{Ipac2}} & Accrued Interest 2 & The current value of accrued interest of the second leg \\ \hline
	\hypertarget{state:Ipnr}{\texttt{Ipnr}} & Nominal Interest Rate & The applicable nominal rate \\ \hline
	\hypertarget{state:Ipnr2}{\texttt{Ipnr2}} & Nominal Interest Rate 2 & The applicable nominal rate \\ \hline
	\hypertarget{state:Md}{\texttt{Md}} & Maturity Date & The timestamp as per which the contract matures according to the initial terms or as per unscheduled events \\ \hline
	\hypertarget{state:Npd}{\texttt{Npd}} & Non Performing Date & The date of the (uncovered) payment event responsible for the current value of the Contract Performance state variable. \\ \hline
	\hypertarget{state:Nt}{\texttt{Nt}} & Notional Principal & The outstanding nominal value \\ \hline
	\hypertarget{state:Nt2}{\texttt{Nt2}} & Notional Principal 2 & The outstanding nominal value of the second leg \\ \hline
	\hypertarget{state:Prf}{\texttt{Prf}} & Contract Performance & Indicates the current Contract Performance status. Different states of the contract range from performing to default. \\ \hline
	\hypertarget{state:Prnxt}{\texttt{Prnxt}} & Next Principal Redemption Payment & The value at which principal is being repaid. This may be including or excluding of interest depending on the Contract Type \\ \hline
	\hypertarget{state:Scip}{\texttt{Scip}} & Interest Scaling Multiplier & The multiplier being applied to interest cash flows \\ \hline
	\hypertarget{state:Scnt}{\texttt{Scnt}} & Notional Scaling Multiplier & The multiplier being applied to principal cash flows \\ \hline
	\hypertarget{state:Sd}{\texttt{Sd}} & Status Date & The timestamp as per which the state is captured at any point in time \\ \hline
	\hypertarget{state:Td}{\texttt{Td}} & Termination Date & The timestamp of unscheduled termination of a contract \\ \hline
	\hypertarget{state:Xa}{\texttt{Xa}} & Exercise Amount & The amount fixed for a the contingent event/obligation exercised at Exercise Date \\ \hline
	\hypertarget{state:Xd}{\texttt{Xd}} & Exercise Date & The timestamp at which a contingent event/obligation is exercised \\ \hline
\end{dicttable}

\subsection{Contract Types}

\begin{dicttable}
	\hypertarget{ct:ANN}{\texttt{ANN}} & Annuity & Principal payment fully at IED and interest plus principal repaid periodically in constant amounts till MD. If variable rate, total amount for interest and principal is recalculated to be fully matured at MD. \\ \hline
	\hypertarget{ct:ANX}{\texttt{ANX}} & Exotic Annuity & Exotic version of ANN However step ups with respect to (i) Principal, (ii) Interest rates are possible. Highly flexible to match totally irregular principal payments. Principal can also be paid out in steps. \\ \hline
	\hypertarget{ct:BCS}{\texttt{BCS}} & Boundary Controlled Switch & A boundary controlled switch is a derivative contract with subcontract legs which can be activated (knocked in) or extinguished (knocked out) when the underlying asset price reaches a specified value. The underlying asset may be a stock, index, or exchange-traded fund. Boundary controlled switch contracts with a single boundary are currently defined. \\ \hline
	\hypertarget{ct:BNDCP}{\texttt{BNDCP}} & Convertible Note & Bonds with a call or put option. If option is exercised, underlying bond ceases to exist. \\ \hline
	\hypertarget{ct:BNDWR}{\texttt{BNDWR}} & Warrant & Bonds with a warrant. If option is exercised, underlying bond continues to exist. \\ \hline
	\hypertarget{ct:CAPFL}{\texttt{CAPFL}} & Cap Floors & Interest rate option expressed in a maximum or minimum interest rate. \\ \hline
	\hypertarget{ct:CDSWP}{\texttt{CDSWP}} & Credit Default Swap & A payment is triggered if one or more of the ndelying counterparties defaults. \\ \hline
	\hypertarget{ct:CEC}{\texttt{CEC}} & Collateral & Collateral creates a relationship between a collateral an obligee and a debtor, covering the exposure from the debtor with the collateral. \\ \hline
	\hypertarget{ct:CEG}{\texttt{CEG}} & Guarantee & Guarantee creates a relationship between a guarantor, an obligee and a debtor, moving the exposure from the debtor to the guarantor. \\ \hline
	\hypertarget{ct:CLM}{\texttt{CLM}} & Call Money & Loans that are rolled over as long as they are not called. Once called it has to be paid back after the stipulated notice period. \\ \hline
	\hypertarget{ct:CLNTE}{\texttt{CLNTE}} & Credit Linked Note & A credit linked note is a security with an embedded CDSWP \\ \hline
	\hypertarget{ct:COM}{\texttt{COM}} & Commodity & This is not a financial contract in its propper sense. However it traks movements of commodities such as oil, gas or even houses. Such commodities can serve as underlyings of commodity futures, guarantees or simply asset positions. \\ \hline
	\hypertarget{ct:CSH}{\texttt{CSH}} & Cash & Cash or cash equivalent position \\ \hline
	\hypertarget{ct:EXOTi}{\texttt{EXOTi}} & EXOTi & As of current, most of the exotic options which were popular before 2008 are out of favor and factually irrelevant.Which of the exotic options will be implemented will depend on the real relevance in the future. \\ \hline
	\hypertarget{ct:FUTUR}{\texttt{FUTUR}} & Future & Keeps track of value changes for any basic CT as underlying (PAM, ANN etc. but also FXOUT, STK, COM). Handles margining calls. \\ \hline
	\hypertarget{ct:FXOUT}{\texttt{FXOUT}} & Foreign Ex-change Outright & Two parties agree to exchange two fixed cash flows in different currencies at a certain point in time in future. \\ \hline
	\hypertarget{ct:LAM}{\texttt{LAM}} & Linear Amortizer & Principal payment fully at IED. Principal repaid periodically in constant amounts till MD. Interest gets reduced accordingly. If variable rate, only interest payment is recalculated. Fixed and variable rates. \\ \hline
	\hypertarget{ct:LAX}{\texttt{LAX}} & Exotic Linear Amortizer & Exotic version of LAM. However step ups with respect to (i) Principal, (ii) Interest rates are possible. Highly flexible to match totally irregular principal payments. Principal can also be paid out in steps. \\ \hline
	\hypertarget{ct:MAR}{\texttt{MAR}} & Margining & A Margining contract traces the value changes and the different margin categories like inital and variation margin. \\ \hline
	\hypertarget{ct:NAM}{\texttt{NAM}} & Negative Amortizer & Similar as ANN. However when resetting rate, total amount (interest plus principal) stay constant. MD shifts. Only variable rates. \\ \hline
	\hypertarget{ct:NAX}{\texttt{NAX}} & Exotic Negative Amortizer & Exotic version of NAM However step ups with respect to (i) Principal, (ii) Interest rates are possible. Highly flexible to match totally irregular principal payments. Principal can also be paid out in steps. \\ \hline
	\hypertarget{ct:OPTNS}{\texttt{OPTNS}} & Option & Calculates straight option pay-off for any basic CT as underlying (PAM, ANN etc.) but also SWAPS, FXOUT, STK and COM. Single, periodic and continuous strike is supported. \\ \hline
	\hypertarget{ct:PAM}{\texttt{PAM}} & Principal at Maturity & Principal payment fully at Initial Exchange Date (IED) and repaid at Maturity Date (MD). Fixed and variable rates. \\ \hline
	\hypertarget{ct:PBN}{\texttt{PBN}} & Perpetual Bonds & Bonds without any maturity date. Interest is paid into eternity if is not terminated. \\ \hline
	\hypertarget{ct:REP}{\texttt{REP}} & Repurchase Agreement & A Repo contract controls and manages the sale and repurchase of assets on the books. \\ \hline
	\hypertarget{ct:SCRCR}{\texttt{SCRCR}} & Securitization Credit Risk & Securitiazion contracts where contracs are ranked according to credit default. The lower ranked tranches are hit by the first defaults. Only when the lowest tranches are wiped out, the next higher tranch is hit. \\ \hline
	\hypertarget{ct:SCRMR}{\texttt{SCRMR}} & Securitization Market Risk & Securitiazion contracts where all underlying contracts are treated equal. The buyer of a tranche gets a part of the cash-flows. \\ \hline
	\hypertarget{ct:STK}{\texttt{STK}} & Stock & Any instrument which is bought at a certain amount (market price normally) and then follows an index. \\ \hline
	\hypertarget{ct:SWAPS}{\texttt{SWAPS}} & Swap & Exchange of two basic CT´s (PAM, ANN etc.). Normally one is fixed, the other variable. However all variants possible including different currencies for cross currency swaps, basic swaps or even different principal exchange programs. \\ \hline
	\hypertarget{ct:SWPPV}{\texttt{SWPPV}} & Plain Vanilla Swap & Plain vanilla swaps where the underlyings are always two identical PAM´s however with one leg fixed and the other variable. \\ \hline
	\hypertarget{ct:TRSWP}{\texttt{TRSWP}} & Total Return Swap & A total return swap is a swap agreement in which one party makes payments based on a set rate, either fixed or variable, while the other party makes payments based on the return of an underlying asset, which includes both the income it generates and any capital gains. \\ \hline
	\hypertarget{ct:UMP}{\texttt{UMP}} & Undefined Maturity Profile & Principal paid in and out at any point in time without prefixed schedule. Interest calculated on outstanding and capitalized periodically. Needs link to a behavioral function describing expected flows. \\ \hline
\end{dicttable}
